\documentclass[uplatex,a4paper,12pt]{jsarticle}
\usepackage{amsmath,amssymb}
\usepackage{enumitem}

\begin{document}

\title{ニューラルネットワークの最適化・正則化 確認テスト}
\author{}
\date{}
\maketitle

\section*{注意}
\begin{itemize}
\item 試験時間は5分とする．
\item 全10問，各10点，合計100点とする．
\item 60点以上を合格とする．
\end{itemize}

\section*{問題}

\begin{enumerate}

\item SGDの説明として最も適切なものを1つ選べ．
\begin{enumerate}[label=(\alph*)]
\item 常に全データを用いて更新する手法
\item 一部のサンプルを用いてパラメータ更新を行う手法
\item 勾配を用いない最適化手法
\end{enumerate}

\item データ集合を1周学習する単位を何というか．

\item モメンタム法の説明として正しいものを1つ選べ．
\begin{enumerate}[label=(\alph*)]
\item 前回の更新量を利用する
\item 学習率を常に0にする
\item 重みを固定する
\end{enumerate}

\item Adamの特徴として正しいものを1つ選べ．
\begin{enumerate}[label=(\alph*)]
\item RMSPropを発展させた手法
\item 勾配を全く利用しない
\item 正則化専用の手法
\end{enumerate}

\item Xavier初期化が主に想定している活性化関数を1つ選べ．
\begin{enumerate}[label=(\alph*)]
\item ReLU
\item sigmoidやtanh
\item 恒等関数
\end{enumerate}

\item He初期化が主に想定している活性化関数を1つ選べ．
\begin{enumerate}[label=(\alph*)]
\item sigmoid
\item tanh
\item ReLU
\end{enumerate}

\item バッチ正規化の目的として最も適切なものを1つ選べ．
\begin{enumerate}[label=(\alph*)]
\item 内部共変量シフトを抑える
\item データ数を増やす
\item クラス数を削減する
\end{enumerate}

\item L2正則化の別名を答えよ．

\item 次のうち，過学習を抑える手法として最も適切なものを1つ選べ．
\begin{enumerate}[label=(\alph*)]
\item Dropout
\item 行列積
\item ソフトマックス関数
\end{enumerate}

\item 転移学習の説明として正しいものを1つ選べ．
\begin{enumerate}[label=(\alph*)]
\item 他タスクで学習した知識を利用する
\item 必ず最初から学習する
\item 教師データを使わない
\end{enumerate}

\end{enumerate}

\end{document}